\section{Quadratic density zeros}
\label{sec:quadratic-density-nonuniqueness}

Infinite-order vanishing of the density is not needed for smooth nonuniqueness. We change the coefficients of $\omega_0$ so that its coefficient of $R$ stays positive and one of its coefficients of $P$ and $Q$ vanishes quadratically at each of the two hypersurfaces $M_\pm$. The resulting singular differential equation forces the slopes of differences of solutions to decay like $\exp(-A/d)$, where $d$ is the distance to the hypersurface and $A>0$. The slopes therefore extend smoothly by zero, even though the density vanishes only to finite order.

\Needspace{14\baselineskip}
\begin{theorem}\label{thm:quadratic-density-nonuniqueness}\index{Zero set!quadratic vanishing}
On $\mathbb P^3$ there exist a smooth nonclosed Hermitian form $\omega$ and a smooth nonnegative density $g$ such that:
\begin{enumerate}
\item The zero set of $g$ consists of two disjoint smooth real hypersurfaces. Near either hypersurface, $g$ is the square of a smooth defining function times a smooth strictly positive function.
\item The equation $(\omega+\ddc u)^3=g\omega_X^3$ has uncountably many distinct smooth $\omega$-plurisubharmonic solutions normalized by $\sup_Xu=0$.
\end{enumerate}
The background $\omega$ has bounded mass.
\end{theorem}

We keep the rank-one forms $P,Q,R$ and the function $r$ of \cref{subsec:numa-rank-directions}, and use new scalar coefficients.

Choose a smooth function $a$ valued in $[0,1]$, equal to one for $r\leq-2/3$ and zero for $r\geq-1/3$, and put $b(r)=a(-r)$. Choose a smooth even cutoff $\chi$ valued in $[0,1]$, zero for $|r|\leq5/4$ and one for $|r|\geq3/2$. Set $p(r)=(1-r^2)^2$ and
\begin{equation}\label{eq:quadratic-background}
 \begin{split}
 x&=(1-\chi)(p+a)+\chi\sigma,\\
 y&=(1-\chi)(p+b)+\chi(1-\sigma),\\
 z&=1-\chi+\chi\sigma',\qquad
 \omega_0=xP+yQ+zR.
 \end{split}
\end{equation}
Since these coefficients coincide with those of \eqref{eq:numa-fs} for $|r|\geq3/2$, the form $\omega_0$ extends by $\omega_X$ across the two omitted projective lines. The coefficient $z$ is strictly positive, and $x,y$ are strictly positive except that $y(-1)=0$ and $x(1)=0$. Thus $\omega_0$ is positive off $M_\pm=\{r=\pm1\}$ and has rank two there.

Choose a smooth cutoff $\kappa\in C_c^\infty((-2,2))$ equal to one on $[-3/2,3/2]$. For a parameter $\varepsilon>0$ to be chosen, define the correction
\[
 \rho=\varepsilon(r^2-1)\kappa(r),\qquad
 \omega=\omega_0+\ddc\rho.
\]
For sufficiently small $\varepsilon>0$, the form $\omega$ is strictly positive. Indeed, on $[-3/2,3/2]$ its coefficients are $x+2\varepsilon r$, $y-2\varepsilon r$ and $z+2\varepsilon$. The first is positive for $r\geq0$, and on $[-3/2,0]$ the function $x$ has a positive minimum, so it remains positive for small $\varepsilon$. Reflection gives the positivity of the second coefficient, and that of the third is immediate. On $3/2\leq|r|\leq2$, the form $\omega_0$ is strictly positive on a compact set and the correction is $O(\varepsilon)$. Near the projective ends there is no correction. The form $\omega$ is not closed: at $r=0$, $x'=y'=0$ and $z=1$, so by \eqref{eq:numa-differential}
\[
 \mathrm{d}\omega=\mathrm{d}r\wedge\bigl((x'-z)P+(y'+z)Q\bigr)\ne0.
\]

Define the density $g$ by
\[
 g\omega_X^3=\omega_0^3.
\]
It is smooth and positive exactly off $M_-\cup M_+$. By \eqref{eq:numa-fs} and \eqref{eq:numa-cube}, $g=xyz/(\sigma(r)(1-\sigma(r))\sigma'(r))$. Near either hypersurface, $xyz=(1+p)p$, and hence
\[
 g(r)=\frac{(1+p(r))(1-r^2)^2}{\sigma(r)(1-\sigma(r))\sigma'(r)}.
\]
The denominator is smooth and strictly positive at $\pm1$, and $1-r^2$ is a defining function for $M_\pm$ near these hypersurfaces. This proves the asserted quadratic vanishing. The potential $-\rho$ is one smooth solution.

\Needspace{15\baselineskip}
\begin{lemma}\label{lem:quadratic-density-trapping}
For every sufficiently small $\eta>0$, the equation
\begin{equation}\label{eq:quadratic-density-ode}
 h'=\frac{xyz}{(x+h)(y-h)}-z,\qquad h(0)=-\eta
\end{equation}
has a smooth solution on $(-1,1)$ with $-x<h<y$ and $h<0$. There are a constant $0<c<1$ and fixed neighborhoods of both endpoints, independent of $\eta$, on which
\[
 -cp<h<0.
\]
Every such solution, extended by zero, is a smooth function on $\mathbb R$ that vanishes to infinite order at $\pm1$.
\end{lemma}
\begin{proof}
The right-hand side $F(r,h)$ is smooth on $-x<h<y$, and $h=0$ is a solution. Near $-1$, $(x,y,z)=(1+p,p,1)$; near $1$, $(x,y,z)=(p,1+p,1)$. Fix $0<c<1$. For the trial function $-cp$, near the left endpoint
\[
 F(r,-cp)+cp'
 =\frac{1+p}{(1+(1-c)p)(1+c)}-1+cp'
 \longrightarrow-\frac{c}{1+c}<0,
\]
whereas near the right endpoint
\[
 F(r,-cp)+cp'
 =\frac{1+p}{(1-c)(1+(1+c)p)}-1+cp'
 \longrightarrow\frac{c}{1-c}>0.
\]
as $r\to\pm1$, since $p,p'\to0$ at both endpoints. The strict inequalities therefore hold on fixed neighborhoods of the endpoints, and $-cp$ is a lower barrier for the equation solved backwards near $-1$ and forwards near $1$.

Choose a compact interval $[r_L,r_R]\ni0$ with endpoints in those neighborhoods. Continuous dependence on initial values, around the solution $h=0$, gives $\eta_0>0$ such that for every $0<\eta<\eta_0$ the solution exists on $[r_L,r_R]$, stays in the domain of $F$, and lies above $-cp$ at $r_L$ and $r_R$. By uniqueness of solutions, it cannot meet the solution zero, so $h<0$. The barrier argument used in the proof of \cref{lem:numa-thin-trapping} then gives $-cp<h<0$ near both ends. On each compact subinterval of $(-1,1)$ the trapped graph stays at a positive distance from the boundary of the domain of $F$; hence, by the continuation theorem, the solution exists on all of $(-1,1)$.

\Needspace{6\baselineskip}
The equation factors as
\begin{equation}\label{eq:quadratic-density-logarithm}
 h'=h\,L(r,h),\qquad
 L(r,h)=\frac{z(x-y+h)}{(x+h)(y-h)}.
\end{equation}
Since $h<0$, we have $(\log(-h))'=L(r,h)$. Near the left endpoint, the trapping bounds and the explicit coefficients imply
\[
 L(r,h)=\frac{1+h}{(1+p+h)(p-h)}\geq\frac{C_1}{p}
\]
after shrinking that neighborhood, with $C_1>0$. Near the right endpoint they give
\[
 L(r,h)=\frac{-1+h}{(p+h)(1+p-h)}\leq-\frac{C_2}{p},
 \qquad C_2>0.
\]
Indeed, on the left, $1+h\geq1-cp\geq1/2$ and the denominator is at most $(1+p)(1+c)p$; on the right, the numerator is at most $-1$, and the denominator is at most $p(1+(1+c)p)$. Since $p$ is comparable to the square of the distance $d=1-|r|$ to the endpoint, integrating these inequalities backwards on the left and forwards on the right gives, for some $A>0$,
\begin{equation}\label{eq:quadratic-density-exponential}
 |h(r)|\leq C \mathrm{e}^{-A/d}
\end{equation}
near either endpoint. The constants may depend on the solution; uniformity in $\eta$ is not needed for the smooth extension.

It remains to control all derivatives. On the trapped graph, all partial derivatives in $(r,h)$ of the rational function $L$ grow at most polynomially in $d^{-1}$. Indeed, near an endpoint $x,y$ are polynomials in $r$, $z=1$, $h$ is bounded, and of the two factors in the denominator, one is bounded below by a positive constant and the other by $(1-c)p$. Repeated differentiation of the quotient therefore only introduces finite powers of $p^{-1}$.

Inductively, for every $j\geq0$ there are constants $C_j,N_j$ with
\[
 |h^{(j)}(r)|\leq C_j \mathrm{e}^{-A/d}d^{-N_j}.
\]
The case $j=0$ is \eqref{eq:quadratic-density-exponential}. Differentiate $h'=hL(r,h(r))$ $j$ times. Every resulting term is the product of a derivative of the factor $h$, of order at most $j$, and a derivative of the composite function $L(r,h(r))$. The latter involves only partial derivatives of $L$ and derivatives of $h$ of order at most $j$, which are polynomially bounded by the induction hypothesis, while the first factor has the exponential bound. This proves the case $j+1$. All derivatives of $h$ therefore tend to zero at $\pm1$. Extending them by zero and applying the mean value theorem to the difference quotients at the endpoints, as in \cref{lem:numa-thin-flat-extension}, shows that the extension of $h$ by zero is smooth and vanishes to infinite order at $\pm1$.
\end{proof}

\begin{proof}[of \cref{thm:quadratic-density-nonuniqueness}]
Let $\eta_0>0$ be such that \cref{lem:quadratic-density-trapping} applies for $0<\eta<\eta_0$, and let $h_\eta$ denote the solution with $h_\eta(0)=-\eta$. Take a primitive $G_\eta$ of its extension by zero. Since $G_\eta$ is constant near each end of the real line, $G_\eta(r)$ defines a smooth function on $X$. The potential $u_\eta=G_\eta(r)-\rho$ satisfies
\[
 \omega+\ddc u_\eta
 =(x+h_\eta)P+(y-h_\eta)Q+(z+h_\eta')R.
\]
Inside the band, the first two coefficients are positive by the lemma, and the equation gives
\[
 z+h_\eta'=\frac{xyz}{(x+h_\eta)(y-h_\eta)}>0.
\]
Outside the band the form is $\omega_0$, and on the two boundary hypersurfaces it is also the rank-two form $\omega_0$, since $h_\eta$ vanishes to infinite order there. Thus it is semipositive on $X$, and by \eqref{eq:numa-cube} its cube is $6xyzP\wedge Q\wedge R=\omega_0^3=g\omega_X^3$. Different initial values give potentials that do not differ by constants, so after normalization by the supremum we obtain an uncountable family. Bounded mass follows from \cref{prop:bmp-metric}, as for \cref{thm:degenerate-hermitian-nonuniqueness}.
\end{proof}

At each of the hypersurfaces $M_\pm$, the difference of two of these solutions agrees to infinite order with a constant, because its derivative is a difference of slopes that vanish to infinite order there. The constants at the two hypersurfaces need not agree. Thus the singular differential equation produces differences that are flat at $M_\pm$, even though the density has only quadratic zeros.
